\documentclass[11pt,a4paper]{article}

\usepackage[utf8]{inputenc}
\usepackage[T1]{fontenc}
\usepackage{geometry}
\geometry{a4paper,left=25mm,right=25mm,top=25mm,bottom=25mm}
\usepackage{amsmath,amssymb}
\usepackage{graphicx}
\usepackage{xcolor}
\usepackage{longtable}
\usepackage{enumitem}
\usepackage{algorithm}
\usepackage{algpseudocode}
\usepackage{hyperref}
\hypersetup{colorlinks,linkcolor=blue,citecolor=black,urlcolor=blue,linktoc=all}
\usepackage{parskip}

\newcommand{\vb}[1]{\boldsymbol{#1}}
\newcommand{\op}[1]{\texttt{#1}}
\newcommand{\field}[1]{\texttt{#1}}
\newcommand{\code}[1]{\texttt{#1}}
\newcommand{\quat}[1]{\hat{q}_{#1}}

\title{Grain Segmentation in exaStamp\\
\large A Workflow Reference}
\author{exaStamp / \texttt{src/grains}}
\date{\today}

\begin{document}
\maketitle

\begin{abstract}
This document describes how exaStamp's grain-segmentation pipeline works: from per-atom crystal
orientation (as fitted by Polyhedral Template Matching, PTM \cite{ptm2016}) to a per-atom grain id,
a per-grain mean orientation, and a color suitable for visualization. The implementation follows
OVITO's own \code{GrainSegmentation} modifier (\code{crystalanalysis/modifier/grains}) directly:
Node-Pair-Sampling graph clustering with an automatically selected merge threshold, and
Dijkstra-by-distance orphan-atom adoption — chosen specifically over the simpler
minimum-spanning-tree/majority-vote alternatives OVITO also offers, for robustness against
thermally noisy samples. The pipeline lives in \texttt{src/grains/} and is exercised end to end by
\texttt{data/regression\_new/grains/compute\_grain\_segmentation.msp}, verified directly against
real OVITO output on that same test case.
\end{abstract}

\tableofcontents

\section{Overview}

Given a per-atom crystal structure type and lattice orientation (from PTM), grain segmentation
partitions the crystalline atoms into contiguous \emph{grains} — regions of consistent
orientation — separated by grain boundaries, and optionally extends each grain to absorb nearby
non-crystalline (``Other'') atoms that sit within its own boundary region. The output, per atom, is
a grain id (or none); per grain, an atom count, a mean lattice orientation, and a display color.

The pipeline has two stages, corresponding to two separate operators:

\begin{longtable}{@{}p{0.32\textwidth}p{0.63\textwidth}@{}}
\hline
\textbf{Operator} & \textbf{Role} \\
\hline
\op{compute\_grain\_bond\_misorientation} &
  GPU-compatible. For every chunk-neighbor bond, records the neighbor's identity, real distance,
  and (for a same-structure-type pair) the crystal-symmetry-aware disorientation angle. \\
\op{compute\_grain\_clusters} &
  Host, sequential. Node-Pair-Sampling clustering with auto-threshold selection, dendrogram cut,
  minimum-grain-size dissolution, optional Dijkstra-by-distance orphan adoption, and per-grain
  color assignment. Writes the final \field{grain\_id}/\field{grain\_color} fields. \\
\hline
\caption{The two grain-segmentation operators.}
\label{tab:pipeline}
\end{longtable}

This split mirrors OVITO's own two-phase engine design (\code{GrainSegmentationEngine1}: neighbor
bonds and the dendrogram, expensive; \code{GrainSegmentationEngine2}: threshold cut and coloring,
cheap) for the same reason: bond disorientation is embarrassingly parallel and belongs on the GPU,
while dendrogram construction and orphan adoption are inherently sequential (each step depends on
the previous one) and are not GPU-parallelizable in OVITO's own implementation either.

\section{Input: PTM structure type and orientation}

Both operators consume two per-particle fields already produced upstream by \op{compute\_ptm} +
\op{ptm\_fields}:
\begin{itemize}[nosep]
  \item \field{ptm\_type}: a discrete structure-type code (none, FCC, HCP, BCC, ICO, SC, \ldots),
  \item \field{ptm\_orientation}: the fitted lattice orientation, stored as a rotation matrix.
\end{itemize}
PTM itself fits this orientation via Theobald's QCP (quaternion-characteristic-polynomial) method
together with combinatorial template/graph matching for the atom-to-template correspondence
\cite{ptm2016} — a well-validated, established registration method, not something this pipeline
re-derives.

Only two symmetry families are supported for clustering, matching OVITO's own
\code{GrainSegmentation} scope exactly: \textbf{cubic} (FCC, BCC, SC — the 24-element proper
rotation point group) and \textbf{HCP} (the 12-element ``conventional'' hexagonal group). Icosahedral
structure has no supported symmetry group in OVITO's own grains code either, and is therefore
never merged into a grain by this pipeline.

Ghost-atom copies of these two fields must be explicitly synchronized (\op{ghost\_update\_opt})
before either grain-segmentation operator runs, since bonds routinely cross into a neighboring
rank's ghost territory.

\section{Bond disorientation}
\label{sec:disorientation}

\subsection{Disorientation angle}

For two atoms of the \emph{same} structure type with fitted orientation quaternions $\quat{0}$
and $\quat{1}$ (converted on the fly from the stored rotation matrix, via a numerically robust,
branch-selected matrix-to-quaternion formula), the \emph{disorientation} is the smallest possible
misorientation angle between the two orientations once every crystallographically equivalent
representation of each is accounted for:
\begin{equation}
  \theta = \min_{g \in G} \; \angle\!\left(\quat{0}^{-1}\quat{1}\, g\right),
\end{equation}
where $G$ is the crystal's own point-group symmetry group (24 elements for cubic, 12 for HCP) acting
on the relative orientation quaternion $q_{\text{rel}} = \quat{0}^{-1}\quat{1}$. Concretely: rotate
$q_{\text{rel}}$ into its own \emph{fundamental zone} by finding the symmetry element $g^\star$
maximizing $|q_{\text{rel}}\cdot g|$ (the crystallographically equivalent orientation closest to the
identity), then
\begin{equation}
  \theta = \arccos\!\left(2t^2-1\right), \qquad t = \left(q_{\text{rel}}\, g^\star\right)_0 ,
\end{equation}
the scalar (real) component of the fundamental-zone-rotated relative quaternion — the standard
double-angle identity $\cos\theta = 2\cos^2(\theta/2)-1$ applied to a unit quaternion's own scalar
part, with the squaring removing the quaternion double-cover sign ambiguity. A bond between two
atoms of \emph{different} structure types, or where either endpoint has no structure classification
at all, has no defined disorientation.

This formula, together with the two symmetry-group tables it uses, is ported verbatim from the
same third-party PTM library \cite{ptm2016} OVITO's own \code{GrainSegmentation} code calls
(\code{ptm\_quat.cpp}) — already vendored elsewhere in exaStamp for PTM itself, but duplicated here
as small, explicitly device-annotated inline functions rather than linked, specifically so this
pass can be genuinely GPU-compatible (the vendored library's own functions are plain host C++, not
written for device compilation).

\subsection{GPU-compatible bond recording}

\op{compute\_grain\_bond\_misorientation} visits every chunk-neighbor bond of every owned atom and,
for each one, records the neighbor's global id, the real Euclidean bond length, and — if both atoms
share a supported structure type and their disorientation clears a hard ceiling
(\code{misorientation\_threshold}, 4$^\circ$ by default, matching OVITO's own hardcoded value) —
the disorientation angle itself; otherwise a sentinel value marks the bond as ineligible for
clustering. Every atom's own set of bonds is capped at a small, fixed number of slots (24); when a
new bond arrives after the cap is reached, it \emph{replaces} whichever currently-recorded bond is
farthest away, so the retained set is always the nearest bonds seen so far, not merely the first
ones encountered in traversal order. This matters for two separate downstream consumers: a dropped
near bond can be the only real candidate for clustering, and — more consequentially — it can be the
only short path an orphan atom needs to reach a grain (\S\ref{sec:orphan}).

The recorded bond list therefore serves two distinct downstream purposes at once, exactly as in
OVITO's own engine (which reuses one \code{\_neighborBonds} array for both): the disorientation
value gates which bonds are \emph{crystalline clustering candidates} (\S\ref{sec:nps}), while the
distance value, together with every recorded bond regardless of disorientation validity, feeds
orphan-atom adoption (\S\ref{sec:orphan}).

\section{Node-Pair-Sampling clustering}
\label{sec:nps}

\subsection{The graph}

Every candidate crystalline bond (disorientation below threshold) becomes one edge of an undirected
weighted graph over atoms, with weight
\begin{equation}
  w = \exp\!\left(-\theta^2/3\right),
\end{equation}
a Gaussian decay in the disorientation angle $\theta$ (in degrees): a perfectly aligned bond gets
weight 1, and weight falls off smoothly as misorientation grows toward the 4$^\circ$ ceiling.

\subsection{Reciprocal-nearest-neighbor chain clustering}

Node-Pair-Sampling is a form of hierarchical agglomerative clustering, implemented as a
reciprocal-nearest-neighbor (RNN) chain — the classical technique that avoids re-scanning the whole
graph after every single merge. For a node $a$, its own nearest neighbor is
\begin{equation}
  \operatorname{NN}(a) = \operatorname*{argmin}_{v \,\in\, \operatorname{adj}(a)} \frac{\operatorname{wnode}(v)}{w(a,v)},
\end{equation}
where $\operatorname{wnode}(v)$ is the sum of edge weights incident on $v$ — a node with many strong
edges elsewhere (already part of a well-formed grain) is a \emph{less} attractive nearest neighbor
than one with few, concentrating merges where the local evidence is strongest. Algorithm~\ref{alg:rnn}
sketches the chain-following procedure.

\begin{algorithm}
\caption{Reciprocal-nearest-neighbor chain clustering}
\label{alg:rnn}
\begin{algorithmic}[1]
\While{active nodes remain}
  \State $\mathit{chain} \gets [\,\text{any active node}\,]$
  \While{$\mathit{chain}$ is not empty}
    \State $a \gets \mathit{chain}.\text{pop}()$
    \State $(d,b) \gets \operatorname{NN}(a)$
    \If{$b$ does not exist} \Comment{$a$'s own connected component is fully clustered}
      \State deactivate $a$
    \ElsIf{$\mathit{chain}$ is not empty}
      \State $c \gets \mathit{chain}.\text{pop}()$
      \If{$c = b$} \Comment{$a$ and $b$ are \emph{mutual} nearest neighbors}
        \State merge $a$ and $b$ (contract the smaller-degree node into the larger)
        \State accumulate the absorbed node's own orientation into the survivor (\S\ref{sec:orient})
        \State record a dendrogram entry at distance $d$
        \State $\mathit{chain}.\text{push}(\text{merged node})$
      \Else
        \State $\mathit{chain}.\text{push}(c)$; $\mathit{chain}.\text{push}(a)$; $\mathit{chain}.\text{push}(b)$
      \EndIf
    \Else
      \State $\mathit{chain}.\text{push}(a)$; $\mathit{chain}.\text{push}(b)$
    \EndIf
  \EndWhile
\EndWhile
\end{algorithmic}
\end{algorithm}

Following the chain — rather than always restarting from scratch — guarantees that when two nodes
finally \emph{are} found to be mutual nearest neighbors, merging them immediately is always a valid,
myopically-optimal step; a node whose own nearest neighbor keeps changing as the chain extends
simply gets pushed further down the chain until a genuine mutual pair is found somewhere along it.
This scales far better than repeatedly rescanning every node's own nearest neighbor from scratch
after each individual merge.

\emph{Implementation note.} exaStamp's own adjacency structure is a plain hash map per node
(neighbor $\to$ edge weight), in place of OVITO's own intrusive red-black tree plus preallocated
edge buffer — the same algorithm and the same clustering result, at the cost of a somewhat less
memory-efficient constant factor that is a complete non-issue at ordinary atomic coordination
numbers ($\sim$12--18 neighbors per atom).

\emph{Why Node-Pair-Sampling, not a minimum spanning tree.} OVITO also offers a simpler
minimum-spanning-tree (Kruskal, single-linkage) clustering mode. Single-linkage clustering is prone
to \emph{chaining}: a handful of thermally noisy, low-misorientation bonds can bridge two
genuinely distinct grains into one cluster via a cheap path, even when the bulk of their own atoms
disagree strongly. The RNN-chain contraction above behaves more like average-linkage clustering
(edge weights accumulate as clusters merge, so a merge decision reflects the accumulated evidence of
an entire growing cluster, not one lone bond) — considerably more resistant to this failure mode,
which is precisely why it was chosen here for studying samples with real thermal noise.

\subsection{Orientation accumulation}
\label{sec:orient}

Each atom begins as its own single unit quaternion. At every merge, the absorbed cluster's own
running orientation sum is first \emph{mapped} onto the surviving cluster's current orientation
neighborhood — the same fundamental-zone search used for the disorientation formula itself, applied
here to find which crystallographically equivalent representation of the absorbed orientation is
actually closest to the survivor — and then accumulated, weighted by its own pre-mapping norm:
\begin{equation}
  \vb{q}_{\text{parent}} \;\mathrel{+}= \; \operatorname{map}(\vb{q}_{\text{parent}}, \vb{q}_{\text{child}}) \cdot \lVert \vb{q}_{\text{child}} \rVert .
\end{equation}
Since every atom starts as a unit quaternion and quaternions are never renormalized between merges,
the norm of a cluster's own running sum equals exactly the number of atoms merged into it so far —
a compact way to implement a symmetry-aware running weighted quaternion mean without ever
re-visiting member atoms. The final per-grain orientation reported downstream is this running sum,
normalized once at the very end.

\section{Automatic threshold selection}
\label{sec:threshold}

The RNN-chain clustering produces a full dendrogram (every merge, in the order it happened,
together with the distance $d$ at which it occurred and the size of each side being merged); the
question of \emph{where} to cut it — which merges represent genuine intra-grain noise and which
represent crossing a real grain boundary — is answered by a robust statistical fit, not a
user-chosen fixed threshold, following OVITO's own default (\code{GraphClusteringAutomatic}) mode
exactly.

Within a single, well-formed grain, successive merges' own distances tend to grow smoothly and
predictably with the size of the clusters being merged (more atoms already agreeing makes further
agreement less surprising); a merge that instead crosses a real grain boundary shows up as an
\emph{anomalous jump} away from that trend. Concretely: fit a power law
$\log d \approx \gamma \log s + c$ (where $s$ is the harmonic-mean merge size) via robust,
iteratively reweighted least-absolute-deviations regression (100 fixed iterations, weight update
$w \gets s / \max(10^{-4}, |\text{residual}|)$ at each iteration); the threshold is then the
\emph{largest} $\log d$ among merges still consistent with that fitted trend (residual below
$1.5\times$ the median absolute residual) — i.e.\ the last merge that still looks like ordinary
intra-grain noise. Every merge above that threshold is presumed to be a genuine inter-grain
boundary crossing and is excluded when the dendrogram is cut.

A manual threshold (OVITO's \code{GraphClusteringManual} mode) is also exposed, expressed in this
same internal log-distance unit rather than a physically interpretable angle — this is an
intentional fidelity choice, matching how OVITO's own manual mode for this particular clustering
algorithm works, not a simplification introduced here.

\section{Cutting the dendrogram}

Applying the threshold is a single linear pass over the (distance-sorted) dendrogram, replaying each
accepted merge through a union-find structure until the first merge whose own $\log d$ exceeds the
threshold. Every resulting connected component below \field{min\_grain\_atom\_count} (100 atoms by
default, matching OVITO's own default) is dissolved back to ``no grain'' rather than kept as a
spuriously tiny cluster. Surviving components are assigned contiguous, 1-based grain ids, ordered
largest-to-smallest by atom count, and each is assigned a random display color drawn from a
fixed-seed HSV distribution (seed 1, matching OVITO's own hardcoded seed exactly, so repeated runs
give reproducible, comparable colors).

\section{Orphan atom adoption}
\label{sec:orphan}

An atom with no PTM structure classification at all (sitting inside a defect core or a genuinely
disordered region) — or a crystalline atom whose own cluster was dissolved for being too small — has
no grain id after the steps above. Orphan adoption (optional, on by default) assigns each such atom
to whichever grain is geometrically closest, using every recorded bond (\S\ref{sec:disorientation}),
not only the crystalline-candidate subset.

This is a single-pass, multi-source Dijkstra: a min-priority queue is seeded with every bond
connecting an already-assigned atom to an unassigned neighbor, weighted by real bond length; the
queue is repeatedly popped, and whichever unassigned atom is reached by the cheapest accumulated
path is settled to that path's own originating grain and used to extend the search further, exactly
as in an ordinary shortest-path computation with multiple sources. An atom is therefore adopted by
the grain that is closest by cumulative real-space path length through the network of
non-crystalline/boundary bonds — not by a majority vote among its own immediate neighbors, and not
by straight-line distance.

\emph{A real bug found via a boundary-region visual inspection, and its fix.} Because bond recording
(\S\ref{sec:disorientation}) only ever runs on owned (not ghost) atoms, a ghost atom's own ``outgoing''
bond list is always empty — even when that ghost carries a perfectly valid grain id (as the
periodic image of an already-clustered atom). Seeding and propagating the priority queue directly
from each atom's own one-directional recorded bond list therefore could never adopt an orphan whose
only nearby assigned neighbor happened to be a ghost, i.e.\ exactly the atoms sitting right at a
periodic domain boundary. The fix builds one \emph{symmetric} adjacency structure from every
recorded bond (both directions, regardless of which atom happened to record it) before running
Dijkstra, removing the directional gap entirely. On the staged 28k-atom, 8-grain test case, this
raised orphan coverage from $\sim$90\% (2913 atoms permanently stuck at boundaries, median distance
to the nearest domain face 0.44~\AA{} versus 4.67~\AA{} over all atoms) to a full 100\%.

\section{Output}

\op{compute\_grain\_clusters} writes two per-atom fields: \field{grain\_id} (integer-valued,
$-1$ for ``no grain'', matching this project's own established convention for a similar field in
the dislocation-extraction pipeline) and \field{grain\_color} (RGB, looked up from the owning
grain's own assigned color, or a fixed neutral gray for unassigned atoms). Both are ordinary named
grid fields and can be written by \op{write\_xyz} or visualized directly in ParaView like any other
per-particle field.

\section{Verification against real OVITO output}

The staged regression test (\texttt{data/regression\_new/grains/}) is a synthetic 28{,}498-atom Ta
polycrystal built from 8 randomly oriented grains in an 80~\AA{} cubic box. Run with
\field{min\_grain\_atom\_count}$=100$ and orphan adoption enabled, and compared directly against a
real OVITO run on the identical input with the same parameters:

\begin{itemize}[nosep]
  \item \textbf{8/8 grains} — exact match to the true grain count.
  \item \textbf{100\% agreement on assigned-vs-unassigned}: all 28{,}498 atoms are either
    grain-assigned in both outputs or in neither.
  \item Per-grain sizes agree with OVITO's own to within \textbf{1--6\%}
    (this implementation: 5539/4574/3818/4080/3988/3554/1560/1385 atoms;
    OVITO: 5494/4561/4066/4021/3827/3596/1547/1386) — a close match given the intentionally
    simplified (not byte-identical) adjacency data structure used for clustering.
\end{itemize}

\section{Known limitations}

\begin{itemize}[nosep]
  \item \textbf{Coherent-interface handling.} OVITO's own optional pre-clustering pass that
    reclassifies a coherent twin boundary or stacking fault (locally a different but closely
    related structure, e.g.\ HCP inside an FCC matrix) into the surrounding grain's own phase
    before clustering is not implemented here; every such interface is currently treated as an
    ordinary grain boundary.
  \item \textbf{Cross-rank MPI grain stitching.} Each MPI rank currently clusters its own local
    (owned + ghost) view of the system independently. A grain whose own true extent spans an MPI
    rank boundary is \emph{not} reconciled into one consistent id across ranks — the same scope
    gap already documented for this project's dislocation-extraction pipeline's own per-rank
    cluster-building stage. Not yet tested at more than one MPI rank at all.
\end{itemize}

\begin{thebibliography}{9}

\bibitem{ptm2016}
P.~M.~Larsen, S.~Schmidt, J.~Schi\o{}tz,
\emph{Robust structural identification via polyhedral template matching},
Modelling and Simulation in Materials Science and Engineering \textbf{24} (2016) 055007.

\end{thebibliography}

\end{document}
